\clearpage
\section{TX-side Gain Flatness}
\label{sec:tx-flatness}
\raggedbottom% no stretched gaps between text and figures (as in the NF section)
% Generated by FR3_Characterization/plot_datasheet_flatness.py -- do not edit by hand.
\newcommand{\flatTxUpperBypassMedian}{1.55}
\newcommand{\flatTxUpperBypassPninety}{2.79}
\newcommand{\flatTxUpperBypassMin}{0.20}
\newcommand{\flatTxUpperBypassMax}{3.63}
\newcommand{\flatTxUpperOptimalMedian}{1.51}
\newcommand{\flatTxUpperOptimalPninety}{2.55}
\newcommand{\flatTxUpperOptimalMin}{0.17}
\newcommand{\flatTxUpperOptimalMax}{3.29}
\newcommand{\flatTxLowerBypassMedian}{1.21}
\newcommand{\flatTxLowerBypassPninety}{2.22}
\newcommand{\flatTxLowerBypassMin}{0.19}
\newcommand{\flatTxLowerBypassMax}{2.74}
\newcommand{\flatTxLowerOptimalMedian}{1.26}
\newcommand{\flatTxLowerOptimalPninety}{2.19}
\newcommand{\flatTxLowerOptimalMin}{0.16}
\newcommand{\flatTxLowerOptimalMax}{3.07}
\newcommand{\flatRxLowBypassMedian}{1.03}
\newcommand{\flatRxLowBypassPninety}{1.76}
\newcommand{\flatRxLowBypassMin}{0.46}
\newcommand{\flatRxLowBypassMax}{2.63}
\newcommand{\flatRxLowOptimalMedian}{0.95}
\newcommand{\flatRxLowOptimalPninety}{1.48}
\newcommand{\flatRxLowOptimalMin}{0.39}
\newcommand{\flatRxLowOptimalMax}{2.24}
\newcommand{\flatRxHighBypassMedian}{0.95}
\newcommand{\flatRxHighBypassPninety}{1.74}
\newcommand{\flatRxHighBypassMin}{0.37}
\newcommand{\flatRxHighBypassMax}{2.89}
\newcommand{\flatRxHighOptimalMedian}{0.90}
\newcommand{\flatRxHighOptimalPninety}{1.81}
\newcommand{\flatRxHighOptimalMin}{0.38}
\newcommand{\flatRxHighOptimalMax}{3.19}
\newcommand{\flatTxPpRatio}{3.3}
\newcommand{\flatTxPpRatioLow}{2.8}
\newcommand{\flatTxPpRatioHigh}{4.1}
\newcommand{\flatRxPpRatio}{4.5}
\newcommand{\flatRxPpRatioLow}{3.7}
\newcommand{\flatRxPpRatioHigh}{5.2}


\textbf{What is gain flatness?} Sections~\ref{sec:tx-gain-upper} and \ref{sec:tx-gain-lower} give the
\emph{average} conversion gain across a 500\,MHz window. A wideband signal also cares how much the gain varies
\emph{within} that window: if one part of the channel is amplified more than another, the signal comes out with a
tilted or rippled spectrum, and the pieces of it that sit at different frequencies (for example the sub-carriers of
an OFDM signal) arrive at different levels. A perfectly flat channel has the same gain at every frequency in the
window. Gain flatness measures how far the unit is from that ideal.

\textbf{How it's measured.} At every (FR1, FR3) point of the grid, the FieldFox's tracking generator sweeps a
500\,MHz window centred on FR1 into the ``TX-side FR1 IF In'' port at −40\,dBm, in 101 steps 5\,MHz apart, while
the FieldFox's receiver follows the output at the ``TX-side FR3 RF Out'' port in lock-step. This gives the
conversion gain at 101 frequencies across the window. The signal level arriving at the unit was measured directly
beforehand, by connecting the FieldFox's source to its own receiver through the same input cable, so any unevenness
of the instrument and the cable is removed instead of being counted against the unit; the output cable's loss is
removed from its own calibration. Each window is measured twice, back to back, with the LO unchanged: once with the
AM3153 filter bypassed and once with the ``optimal'' filter, chosen by the same rule as everywhere else in this
datasheet (Section~\ref{sec:tx-gain-upper}). The filter setting never changes during a sweep. The grid is the same as
for conversion gain: FR1 centred at 1.0--4.0\,GHz and FR3 at 6.0--26.0\,GHz, both in 0.5\,GHz steps, wherever the
LO lands between 7 and 22.6\,GHz.\footnote{The one grid point at FR3 = 26.5\,GHz (upper sideband, FR1 = 4.0\,GHz) is
not measured: its window would run past the 26.5\,GHz top of the FieldFox's range.}

\textbf{Measurement conditions.} \lmxDriveRule{} (set from the LO frequency), with the LO leakage suppressed by the unit's stored LO-null calibration (Section~\ref{sec:lo-suppression}), exactly as in normal use. The unit was measured hot, as the worst case: no fan, LMX2820 die at \condFlatMinC--\condFlatMaxC\textdegree C. A cooler unit will typically do somewhat better.

\textbf{The number: RMS flatness.} Each window gives 101 gain readings G\textsubscript{1} to G\textsubscript{101}, in
dB. First, the window's power-average gain is found: average the gains as linear power ratios, then convert back to
dB,
\begin{center}
G\textsubscript{avg} = 10·log\textsubscript{10}( mean of 10\textsuperscript{G\textsubscript{k}/10} ),
\end{center}
which is the single, constant gain that would deliver the same total power (the window average plotted in
Sections~\ref{sec:tx-gain-upper} and \ref{sec:tx-gain-lower}). The RMS flatness is then the root-mean-square of
how far each reading sits from that level:
\begin{center}
RMS flatness = √( mean of (G\textsubscript{k} − G\textsubscript{avg})\textsuperscript{2} ).
\end{center}
0\,dB means perfectly flat; the larger the number, the more uneven the gain across the window. For a feel of
scale: in this data the full spread from the lowest to the highest gain in a window (peak-to-peak) is typically
about \flatTxPpRatio{} times the RMS flatness (between \flatTxPpRatioLow{} and \flatTxPpRatioHigh{} times at 80\% of
the points).

\textbf{Results.} Within the design spec (FR3~$\le$~24\,GHz, and FR1 centre~$\le$~3.5\,GHz, since the window at
FR1 = 4.0\,GHz runs past the IF spec):
\par\medskip\noindent\begin{minipage}{\linewidth}% caption and table stay on one page
\begin{center}
\captionof{table}{TX-side RMS gain flatness within the design spec, by sideband and filter setting.}\label{tab:txflat-results}
\vspace{0.4em}
\begin{tabular}{lcccc}
\toprule
 & Median & 90\% of points at or below & Best & Worst \\
\midrule
Upper sideband, filter bypassed & \flatTxUpperBypassMedian\,dB & \flatTxUpperBypassPninety\,dB & \flatTxUpperBypassMin\,dB & \flatTxUpperBypassMax\,dB \\
Upper sideband, optimal filter & \flatTxUpperOptimalMedian\,dB & \flatTxUpperOptimalPninety\,dB & \flatTxUpperOptimalMin\,dB & \flatTxUpperOptimalMax\,dB \\
Lower sideband, filter bypassed & \flatTxLowerBypassMedian\,dB & \flatTxLowerBypassPninety\,dB & \flatTxLowerBypassMin\,dB & \flatTxLowerBypassMax\,dB \\
Lower sideband, optimal filter & \flatTxLowerOptimalMedian\,dB & \flatTxLowerOptimalPninety\,dB & \flatTxLowerOptimalMin\,dB & \flatTxLowerOptimalMax\,dB \\
\bottomrule
\end{tabular}
\end{center}
\end{minipage}\par
Engaging the optimal filter changes the flatness only a little: its median is slightly lower than with the filter
bypassed, for both sidebands.

\textbf{Where it's flattest, and how to use that.} The flatness depends mostly on where the signal sits in FR1, and
much less on FR3 or on the sideband. In the heatmaps this shows up as whole columns of similar colour. The median per
FR1 centre (all FR3 up to 24\,GHz):
\par\medskip\noindent\begin{minipage}{\linewidth}% caption and table stay on one page
\begin{center}
\captionof{table}{TX-side median RMS gain flatness for each FR1 centre (all FR3 up to 24\,GHz).}\label{tab:txflat-by-fr1}
\vspace{0.4em}
\footnotesize
% Generated by FR3_Characterization/plot_datasheet_flatness.py -- do not edit by hand.
\begin{tabular}{lcccc}
\toprule
FR1 centre & Upper sideband, bypassed & Upper sideband, optimal & Lower sideband, bypassed & Lower sideband, optimal \\
\midrule
1.0\,GHz & 0.55\,dB & 0.42\,dB & 0.53\,dB & 0.43\,dB \\
1.5\,GHz & 2.07\,dB & 1.70\,dB & 2.17\,dB & 1.89\,dB \\
2.0\,GHz & 2.17\,dB & 1.92\,dB & 1.62\,dB & 1.33\,dB \\
2.5\,GHz & 1.28\,dB & 1.50\,dB & 1.06\,dB & 1.30\,dB \\
3.0\,GHz & 0.65\,dB & 0.72\,dB & 0.70\,dB & 0.82\,dB \\
3.5\,GHz & 2.42\,dB & 2.24\,dB & 1.85\,dB & 2.13\,dB \\
4.0\,GHz\textsuperscript{*} & 2.01\,dB & 1.18\,dB & 1.23\,dB & 1.17\,dB \\
\bottomrule
\multicolumn{5}{l}{\footnotesize\textsuperscript{*}Its 500\,MHz window (3.75--4.25\,GHz) runs 250\,MHz past the IF design spec.}
\end{tabular}

\end{center}
\end{minipage}\par
Windows centred at FR1 = 1.0 or 3.0\,GHz are the flattest, and windows centred at 1.5--2.0 or 3.5\,GHz the least
flat. So when flatness matters and you are free to choose the IF, centre the signal near 1.0 or 3.0\,GHz in FR1.

\textbf{What's plotted.} For each sideband and filter setting, a heatmap of the RMS flatness over the full measured
(FR1, FR3) grid, and a histogram of the same values. All flatness heatmaps in this datasheet (TX and RX) share one
colour scale, 0--4\,dB, darker meaning flatter. On the heatmaps, points above the rated 24\,GHz are boxed in red, as
is the FR1 = 4.0\,GHz column, whose window runs 250\,MHz past the IF spec: measured and shown, but out of design
spec. The histograms count only the points within the design spec.

\figpair{images/flat_tx_upper_bypass_map.pdf}{RMS gain flatness, upper sideband: full measured (FR1, FR3) grid,
filter bypassed. Zoom in to read each cell's exact values, printed as (FR1, FR3, RMS flatness), with FR1 and FR3 in
GHz and the RMS flatness in dB.}{fig:txflat-1}%
        {images/flat_tx_upper_bypass_hist.pdf}{Sideband: Upper. Filter: Bypassed. Distribution of
Figure~\ref{fig:txflat-1}, points within the design spec.}{fig:txflat-2}
\figpair{images/flat_tx_upper_optimal_map.pdf}{RMS gain flatness, upper sideband: full measured (FR1, FR3) grid,
``optimal'' filter engaged at every point.}{fig:txflat-3}%
        {images/flat_tx_upper_optimal_hist.pdf}{Sideband: Upper. Filter: Optimal. Distribution of
Figure~\ref{fig:txflat-3}, points within the design spec.}{fig:txflat-4}
\clearpage
\figpair{images/flat_tx_lower_bypass_map.pdf}{RMS gain flatness, lower sideband: full measured (FR1, FR3) grid,
filter bypassed.}{fig:txflat-5}%
        {images/flat_tx_lower_bypass_hist.pdf}{Sideband: Lower. Filter: Bypassed. Distribution of
Figure~\ref{fig:txflat-5}, points within the design spec.}{fig:txflat-6}
\figpair{images/flat_tx_lower_optimal_map.pdf}{RMS gain flatness, lower sideband: full measured (FR1, FR3) grid,
``optimal'' filter engaged at every point.}{fig:txflat-7}%
        {images/flat_tx_lower_optimal_hist.pdf}{Sideband: Lower. Filter: Optimal. Distribution of
Figure~\ref{fig:txflat-7}, points within the design spec.}{fig:txflat-8}
\flushbottom
